\subsection{弗雷德霍姆映射}

设 E 与 F 是局部凸空间。在本节中，我们用 \(u \equiv v\) 表示模 \(\mathcal{L}^{\mathrm{f}}(\mathrm{E};\mathrm{F})\) 的同余关系，这里 u 与 v 是 \(\mathcal{L}(\mathrm{E};\mathrm{F})\) 的元素。

如果 G 是局部凸空间，且 \(u'\) 与 \(v'\) 是 \(\mathcal{L}(\mathrm{F};\mathrm{G})\) 的元素，那么关系 \(u \equiv v\) 与 \(u' \equiv v'\) 蕴含关系 \(u' \circ u \equiv v' \circ v\)。

我们称元素 \(w\)（属于 \(\mathcal{L}(\mathrm{F};\mathrm{E})\)）是元素 \(u\)（属于 \(\mathcal{L}(\mathrm{E};\mathrm{F})\)）的拟逆，如果有 \(w \circ u \equiv 1_{\mathrm{E}}\) 与 \(u \circ w \equiv 1_{\mathrm{F}}\)。

假设 w 是 u 的拟逆。如果 \(u_{1}\) 是 \(\mathcal{L}(\mathrm{E};\mathrm{F})\) 的元素，且 \(w_{1}\) 是 \(\mathcal{L}(\mathrm{F};\mathrm{E})\) 的元素，使得 \(u_{1} \equiv u\) 且 \(w_{1} \equiv w\)，那么 \(w_{1}\) 是 \(u_{1}\) 的拟逆，因为 \(w \circ u \equiv w_{1} \circ u_{1}\) 且 \(u \circ w \equiv u_{1} \circ w_{1}\)。

如果 w 与 \(w_{1}\) 都是 u 的拟逆，那么 \(w_{1} \equiv w\)，因为

\[
w _ {1} = 1 _ {\mathrm{E}} \circ w _ {1} \equiv w \circ u \circ w _ {1} \equiv w \circ 1 _ {\mathrm{F}} = w.
\]

定义 1. --- 设 E 与 F 是局部凸空间。我们称 \(\mathcal{L}(\mathrm{E};\mathrm{F})\) 中的一个元素 u 为弗雷德霍姆映射 \(^{(1)}\)，如果它拥有一个拟逆。从 E 到 E 中的弗雷德霍姆映射称为 E 的弗雷德霍姆自同态。

我们将用 \(\mathcal{F}(\mathrm{E};\mathrm{F})\) 表示从 E 到 F 中的弗雷德霍姆映射全体所成的集，并用 \(\mathcal{F}(\mathrm{E})\) 表示 E 的弗雷德霍姆自同态全体所成的集。

注记. --- 设 E、F 与 G 是局部凸空间，\(u: E \to F\) 与 \(v: F \to G\) 是连续线性映射。

\begin{enumerate}
\def\labelenumi{\arabic{enumi})}
\item
  假设 u 是弗雷德霍姆映射，并用 \(u_{1}\) 表示 u 的一个拟逆。由于 u 是 \(u_{1}\) 的拟逆，映射 \(u_{1}\) 是弗雷德霍姆映射。
\item
  假设 u 与 v 是弗雷德霍姆映射，并设 \(u_{1}\)（相应地，\(v_{1}\)）是 u（相应地，v）的拟逆。那么 \(v \circ u\) 是从 E 到 G 中的弗雷德霍姆映射，且 \(u_{1} \circ v_{1}\) 是 \(v \circ u\) 的拟逆。事实上，我们计算出
\end{enumerate}

\[
\left(u _ {1} \circ v _ {1}\right) \circ (v \circ u) = u _ {1} \circ \left(v _ {1} \circ v\right) \circ u \equiv u _ {1} \circ 1 _ {\mathrm{F}} \circ u = u _ {1} \circ u \equiv 1 _ {\mathrm{E}},
\]

\[
(v \circ u) \circ (u _ {1} \circ v _ {1}) = v \circ (u \circ u _ {1}) \circ v _ {1} \equiv v \circ 1 _ {\mathrm{F}} \circ v _ {1} = v \circ v _ {1} \equiv 1 _ {\mathrm{G}}.
\]

\begin{enumerate}
\def\labelenumi{\arabic{enumi})}
\setcounter{enumi}{2}
\tightlist
\item
  假设 u 与 \(v \circ u\) 是弗雷德霍姆映射，并设 \(w_{1}\) 是 \(v \circ u\) 的拟逆。那么 v 是弗雷德霍姆映射，且 \(u \circ w_{1}\) 是 v 的拟逆。
\end{enumerate}

事实上，由第一个注记，\(w_{1}\) 是弗雷德霍姆算子。设 \(u_{1}\) 是 u 的拟逆；由第二个注记，\(u \circ w_{1}\) 是弗雷德霍姆映射，且 \((v \circ u) \circ u_{1}\) 是一个拟逆。我们有 \(u \circ u_{1} \equiv 1_{F}\)，从而 \(v \circ u \circ u_{1} \equiv v\)；这就证明了断言。

\begin{enumerate}
\def\labelenumi{\arabic{enumi})}
\setcounter{enumi}{3}
\tightlist
\item
  假设 v 与 \(v \circ u\) 是弗雷德霍姆映射，并设 \(w_{1}\) 是 \(v \circ u\) 的拟逆。那么 u 是弗雷德霍姆算子，且 \(w_{1} \circ v\) 是 u 的拟逆。
\end{enumerate}

证明与前一个注记的证明类似。

引理 1. --- 设 E 与 F 是局部凸空间，u 是从 E 到 F 中的弗雷德霍姆映射。u 的核与余核都是有限维的。

设 \(v \colon \mathrm{F} \to \mathrm{E}\) 是 \(u\) 的拟逆。\(u\) 的核包含于有限秩线性映射 \(1_{\mathrm{E}} - v \circ u\) 的像之中，因而是有限维的。\(u\) 的像包含有限秩线性映射 \(1_{\mathrm{F}} - u \circ v\) 的核，因而在 F 中有有限余维。

命题 2. — 设 E 与 F 是分离的局部凸空间，并设 u 是 \(\mathcal{L}(\mathrm{E};\mathrm{F})\) 的一个元素。下列性质等价：

\begin{enumerate}
\def\labelenumi{(\roman{enumi})}
\item
  映射 u 是弗雷德霍姆算子；
\item
  映射 \(u\) 是严格态射，它的核是有限维的，它的像是闭的且在 F 中有有限余维；
\item
  存在 E 的闭向量子空间 \(E_{1}\) 与 F 的闭向量子空间 \(F_{1}\)，两者都有有限余维，使得 u 通过过渡到子空间诱导了从 \(E_{1}\) 到 \(F_{1}\) 上的同构；
\item
   存在拓扑直和分解 \(E = E_{1} \oplus E_{2}\) 与 \(F = F_{1} \oplus F_{2}\)，使得 \(E_{2}\) 与 \(F_{2}\) 是有限维的，u 在 \(E_{2}\) 上为零，并通过过渡到子空间定义了从 \(E_{1}\) 到 \(F_{1}\) 上的同构。
\item
  \(\Longrightarrow\) (iii)：设 v 是 u 的拟逆，\(E_{1}\) 是 \(1_{E}-v\circ u\) 的核，而 \(F_{1}\) 是 \(1_{F}-u\circ v\) 的核。由于线性映射 \(1_{E}-v\circ u\) 与 \(1_{F}-u\circ v\) 是连续的且是有限秩的，\(E_{1}\) 与 \(F_{1}\) 分别是 E 与 F 的有限余维闭向量子空间。设 \(x\in E_{1}\)。我们有
\end{enumerate}

\[
(1 _ {\mathrm{F}} - u \circ v) (u (x)) = u ((1 _ {\mathrm{E}} - v \circ u) (x)) = u (0) = 0,
\]

从而 \(u(x) \in \mathrm{F}_1\)。因此我们有 \(u(\mathrm{E}_1) \subset \mathrm{F}_1\)；同样我们有 \(v(\mathrm{F}_1) \subset \mathrm{E}_1\)。于是，连续线性映射 \(u_1: \mathrm{E}_1 \to \mathrm{F}_1\) 与 \(v_1: \mathrm{F}_1 \to \mathrm{E}_1\)（它们由 \(u\) 与 \(v\) 得到）是互逆的同构，因为 \(v \circ u\) 与 \(1_{\mathrm{E}}\)（相应地，\(u \circ v\) 与 \(1_{\mathrm{F}}\)）在 \(\mathrm{E}_1\) 上（相应地，在 \(\mathrm{F}_1\) 上）重合。

\begin{enumerate}
\def\labelenumi{(\roman{enumi})}
\setcounter{enumi}{2}
\item
  \(\Longrightarrow\) (ii)：设 \(E_{1}\) 与 \(F_{1}\) 满足 (iii) 的假设。我们有 \(\mathrm{E}_{1}\cap\mathrm{Ker}(u)=\{0\}\) 与 \(\mathrm{F}_{1}\subset\mathrm{Im}(u)\)，因而 \(\mathrm{Ker}(u)\) 是有限维的，而 \(\mathrm{Im}(u)\) 是闭的且在 F 中有有限余维。由 III, p. 39 的 prop. 1 可得映射 u 是严格态射。
\item
   \(\Longrightarrow\) (iv)：假设条件 (ii) 满足。E 的闭向量子空间 \(E_{2} = \operatorname{Ker}(u)\) 是有限维的，且存在 E 的向量子空间 \(E_{1}\)，它在拓扑上与 \(E_{2}\) 互补（EVT, II, p. 27, cor. 2）。向量子空间 \(F_{1} = \operatorname{Im}(u)\) 是闭的且有有限余维，它在 F 中有一个拓扑补 \(F_{2}\)。由 III, p. 39 的 prop. 1，映射 \(u|E_{1}\) 是严格态射，因而 u 诱导了从 \(E_{1}\) 到 \(F_{1}\) 上的同构。
\item
   \(\Longrightarrow\) (i)：在 (iv) 的假设下，从 F 到 E 中的、与 \(u_1^{-1}\) 在 \(\mathrm{F}_1\) 上重合且在 \(\mathrm{F}_2\) 上为零的线性映射是 \(u\) 的拟逆。
\end{enumerate}

注记 5. --- 设 E、F 是分离的局部凸空间，而 \(u\colon\mathrm{E}\to\mathrm{F}\) 是弗雷德霍姆算子。如果 \(u\) 是双射，那么 \(u\) 是同构（事实上，由 Prop. 2, (ii)，\(u\) 是严格态射）。

